\begin{resumo}[Abstract]\begin{otherlanguage*}{english}
Managing a treatment involving multiple medications requires coordinating schedules, quantities, and availability while monitoring daily care events. This work presents DoseCerta, a modular prototype for unit-dose medication dispensing and care monitoring intended for older adults with polypharmacy. The system combines removable reservoirs, a rotating disc with five cavities, a stepper motor, and optical sensing, with distributed control provided by a Raspberry Pi Zero 2 W and ESP32 microcontrollers. The modules communicate over a four-contact interface carrying power and a shared I²C bus. The central unit maintains a local schedule and operational records and synchronizes information with an API and a central database. The Care application provides authenticated configuration, a daily plan, a history, and a user manual, while the computer vision and artificial intelligence agent components support observation and interaction. Development was organized around requirements, incremental construction, and verification of components and interfaces. The evaluation strategy distinguishes software behavior, physical execution, and interpretation of events.
Development records comprise 3,430 communication commands and 12 dispensing cycles, including four classified as retained units, without sufficient independently counted trials to estimate dispensing accuracy. The MCP tool server passed 777 automated tests. The textual evaluation of agent configuration v28 included 210 conversations, with no write request before confirmation in 80 cases and a median latency of 2.3~s per model call. Human review identified time conversion errors, corrected only in a test variant. Scheduled dispensing remained disabled, and final-version vision performance was not measured. The prototype integrates the main system layers and supports interface evaluation; autonomous dispensing and use with people still require physical, timing, and user validation, with no demonstrated clinical efficacy.
\pendencia{PRE-ABS}{Atualizar a tradução caso o resumo seja alterado após novos ensaios; manter versões, denominadores e limitações equivalentes.}
\vspace{\onelineskip}\noindent\textbf{Keywords}: medication dispensing; embedded systems; polypharmacy; systems integration; care support.
\end{otherlanguage*}\end{resumo}
